Bipolar Disorder is routinely mistaken for Major Depressive Disorder, with patients waiting nearly a decade before receiving a correct diagnosis. Both conditions present identically during depressive episodes, yet antidepressants prescribed to unrecognized bipolar patients can precipitate dangerous manic switching. We investigate whether wrist-worn actigraphy can passively distinguish these conditions using the Depresjon dataset (55 participants: 32 healthy controls, 8 bipolar, 15 unipolar). A hybrid 1D-CNN-LSTM architecture operating directly on raw minute-level activity counts is evaluated through four systematic approaches: baseline deep learning with SMOTE, balanced downsampling, four alternative neural architectures, and a diagnostic phase combining classical machine learning with statistical hypothesis testing. Motor activity patterns do not meaningfully differ between bipolar and unipolar depression during depressive episodes. The best reliable accuracy is 63.4\% with 48-hour windows, below the 65.2\% majority-class baseline, and statistical testing confirms negligible group differences (Mann-Whitney $p = 0.776$, Cohen's $d = {-}0.060$). A 5-feature logistic regression matches a 223K-parameter deep network at 60.87\% LOOCV accuracy. These findings constitute a rigorous negative result, ruling out wrist activity variability as a standalone diagnostic biomarker and redirecting future work toward multi-modal and longitudinal approaches.
\keywords{bipolar disorder, unipolar depression, actigraphy, CNN-LSTM, sequence modeling, wearable sensors, differential diagnosis, negative result}
